% =====================================================================
%  Methods — combinatorial perturbations and genetic-interaction scoring
%  Rewritten against sherlock/tools/_models.py (VAE.model, VAE.get_gi,
%  regress_gi_params, classify_gi). Replaces the previous version of
%  \subsection{Combinatorial perturbations and genetic interaction scoring}.
%
%  Cross-references to fix on merge: Sections~4.1--4.4, Eq.~(13),
%  \eqref{eq:full-elbo} are carried over from the old text and must match
%  the surrounding Methods numbering.
%  \ref{sec:causal} points at "Causal Inference of Perturbation Effect";
%  add \label{sec:causal} to that subsection or change the reference.
% =====================================================================

\subsection{Combinatorial perturbations and genetic interaction scoring}
\label{sec:combo}

The sections above describe the single-perturbation case in which each cell carries exactly one
perturbation label $p_n$. We now extend the model to data in which cells may be exposed to a
\emph{pair} of perturbations simultaneously, and introduce a learned interaction term that captures
non-additive effects.

\subsubsection{Notation and input encoding}
\label{sec:combo-notation}

In the combinatorial setting each cell $n$ is associated with an ordered pair
$\mathbf{p}_n = (p_n^{(1)}, p_n^{(2)}) \in \{1,\dots,P\}^2$, where $p_n^{(2)} = \varnothing$
(encoded as $-1$) for single-perturbation cells. Perturbation labels are indexed at the level of
individual genes: a label \texttt{A+B} is split into its components, so a combination introduces no
parameters of its own and instead reuses the embeddings, gates and covariance of its two
constituents. Let
\[
  v_n \;=\; \mathbf{1}\big[p_n^{(2)} \neq \varnothing\big]
\]
indicate that cell $n$ carries a second perturbation. All remaining notation---$A$, $W$, $\rho$,
$z^0$, $z$, $\theta$---is unchanged from Sections~4.1--4.4; in particular $A = \Sigma_P^{1/2}\rho$ is
the covariance-transformed perturbation embedding and $W \in [0,1]^{P\times d}$ the gate.

\subsubsection{Combinatorial latent shift}
\label{sec:combo-shift}

For a single-perturbation cell the gated shift is unchanged,
$\Delta_n = W_{p_n} \odot A_{p_n}$. For a combinatorial cell the shift is a reweighted sum of the two
parent shifts plus a learned interaction residual:
\begin{align}
  \Delta_n^{\mathrm{combo}}
  \;=\;
  c_a\,\big(W_{p_n^{(1)}} \odot A_{p_n^{(1)}}\big)
  \;+\;
  c_b\,\big(W_{p_n^{(2)}} \odot A_{p_n^{(2)}}\big)
  \;+\;
  s\big(\rho_{p_n^{(1)}},\, \rho_{p_n^{(2)}}\big),
  \label{eq:combo-full}
\end{align}
with $c_a = c_b = 1$ and $s = \mathbf{0}$ recovering the additive baseline
$\Delta_n^{\mathrm{add}} = W_{p_n^{(1)}} \odot A_{p_n^{(1)}} + W_{p_n^{(2)}} \odot A_{p_n^{(2)}}$.
Single-perturbation cells bypass \eqref{eq:combo-full} entirely and retain
$\Delta_n = W_{p_n} \odot A_{p_n}$. The shift enters the latent transition exactly as in the
single-perturbation case, with $\Delta_n^{\mathrm{combo}}$ substituted for $\Delta_n$; the decoder,
KL terms and remaining ELBO structure are unchanged.

\paragraph{Interaction residual.}
The residual $s$ is a two-layer perceptron applied to the concatenated pair of perturbation
embeddings and symmetrised over both orderings:
\begin{align}
  s(\rho_a, \rho_b)
  \;=\;
  \tfrac{1}{2}\Big[
    g\big([\rho_a;\rho_b]\big) + g\big([\rho_b;\rho_a]\big)
  \Big],
  \qquad
  g(u) \;=\; W_2\,\mathrm{LeakyReLU}\big(W_1 u + b_1\big) + b_2,
  \label{eq:syn-shift}
\end{align} 
with $W_1 \in \mathbb{R}^{h \times 2d}$, $W_2 \in \mathbb{R}^{d \times h}$ and hidden width $h$.
Averaging the two orderings enforces $s(\rho_a,\rho_b) = s(\rho_b,\rho_a)$, so \texttt{A+B} and
\texttt{B+A} induce identical shifts. The residual is unconstrained in direction and is therefore
free to point outside the span of the two parent shifts, which is what allows it to represent
asymmetric masking (epistasis) and transcriptional states absent from either single perturbation
(neomorphic effects).

\paragraph{Parent coefficients.}
The scalar coefficients $c_a, c_b$ rescale each parent's contribution in the context of the
combination. Both are produced by a single bias-free linear map $w \in \mathbb{R}^{2d}$ applied to
the ordered pair, evaluated in both orderings:
\begin{align}
  c_a \;=\; 1 + \tanh\!\big(w^\top [\rho_a;\rho_b]\big),
  \qquad
  c_b \;=\; 1 + \tanh\!\big(w^\top [\rho_b;\rho_a]\big).
  \label{eq:parent-scale}
\end{align}
This construction has the two properties the interaction classes require. It is swap-equivariant,
$c_a(\rho_a,\rho_b) = c_b(\rho_b,\rho_a)$, giving the label-order invariance appropriate to an
unordered pair. And because the map acts on the full ordered pair, the two coefficients vary
independently and their sum is unconstrained, so a combination can scale its additive component as a
whole---both parents amplified, as in potentiation, or both damped, as in suppression---as well as
redistribute weight between them, as in dominance. The $\tanh$ confines each coefficient to $(0,2)$,
so a parent's program may be damped or up to doubled but cannot change sign or diverge, and $w$ is
initialised near zero so that $c_a = c_b \approx 1$ at initialisation.

\subsubsection{Additivity prior and objective}
\label{sec:combo-elbo}

Nothing in the reconstruction likelihood penalises non-additivity, so without an explicit prior the
model can introduce interaction structure into genuinely additive pairs at negligible cost to the
ELBO---which would erase the additive signature that the genetic-interaction scoring below reads
out. We therefore add two quadratic penalties over combinatorial cells:
\begin{align}
  \mathcal{R}_{\mathrm{syn}}
  \;=\;
  -\lambda_{s} \sum_{n\,:\,v_n=1} \big\| s_n \big\|_2^2
  \;-\;
  \lambda_{c} \sum_{n\,:\,v_n=1} \Big[(c_{a,n}-1)^2 + (c_{b,n}-1)^2\Big].
  \label{eq:syn-reg}
\end{align}
Both accumulate over the cells carrying a combination, as its likelihood contribution does, so how
much evidence a pair must show before departing from additivity does not depend on how many cells
happen to carry it. Together with the near-zero initialisation of $w$, these make the additive model
the default: the model begins exactly at $\Delta^{\mathrm{add}}$ and must be driven away from it by
the data. The complete objective is
\begin{align}
  \mathcal{L}^{\mathrm{combo}}
  \;=\;
  \mathcal{L}\big|_{\Delta \leftarrow \Delta^{\mathrm{combo}}}
  \;+\;
  \mathcal{R}_{\mathrm{syn}},
\end{align}
where $\mathcal{L}$ is \eqref{eq:full-elbo} evaluated with the combinatorial shift
\eqref{eq:combo-full}. The auxiliary classifier loss is the one component that changes form: the
single-label cross-entropy in (13) is replaced by binary cross-entropy against a multi-hot target
$\mathbf{t}_n \in \{0,1\}^P$ with $t_{n,p} = 1$ for each perturbation present in cell $n$,
\begin{align}
  \log p_\psi(y_n \mid z_n)
  \;\equiv\;
  -\mathrm{BCE}\big(\sigma(h_\psi(z_n)),\; \mathbf{t}_n\big),
  \label{eq:combo-ce}
\end{align}
with $\sigma$ the element-wise sigmoid.

\subsubsection{Effect vectors for interaction scoring}
\label{sec:gi-effects}

Interaction scoring operates on effect vectors in gene space---one per single perturbation and one
per pair. In the main analysis each vector is taken from the model's approximate posterior over that
perturbation's own cells: we decode $q(z)$ for every cell carrying the perturbation, average the
decoded profiles, and subtract a non-targeting control (NTC) baseline obtained the same way from the
full NTC population. All three vectors entering a comparison---$\delta_a$, $\delta_b$ and
$\delta_{ab} \in \mathbb{R}^G$---are referenced to this single shared baseline and expressed as
$\log_2$ CPM differences, so a null perturbation maps to $\mathbf{0}$ and the additive residual is
not offset by a baseline mismatch. Perturbations with fewer than five observed cells are not scored.

The alternative is to generate each effect by applying the corresponding latent shift---
\eqref{eq:combo-full} for a pair---to a background abducted from NTC cells and decoding, reusing the
same background throughout. Both estimators target the marginal interventional means identified by
Corollary~\ref{cor:gi-id}, so the choice between them is one of estimator behaviour rather than of
what is identified. We score from the posterior for two reasons. Generating all three vectors makes
the regression close to degenerate by construction: the double is then a reweighted sum of the
singles plus the interaction residual, so the fitted coefficients and residual recover the model's own
parameters rather than a comparison against data. And the reference categories were assigned from
measured joint profiles, so estimating the joint effect from the cells that received the pair mirrors
the construction those labels came from.

For combinations withheld from training we instead generate the pair's effect from
\eqref{eq:combo-full} while retaining posterior-decoded singles, so that the pair is predicted from
its two constituents and its own cells are not used. Regressing a generated double against observed
singles also keeps the interaction residual discriminative.

The status of this setting should be stated precisely, since it differs from the rest of the
analysis. Identification is relative to the distribution a model is fitted to: withholding a
combination removes the very conditional that identifies its mechanism, so the resulting prediction
is an extrapolation rather than an identified counterfactual, obtained by applying parent
coefficients and an interaction residual fitted on other pairs. What makes the extrapolation possible
is structural---a combination reuses its constituents' parameters and therefore has a fully specified
shift whether or not it was observed---and whether it is accurate is an empirical question. We make
no identification claim for withheld combinations. Because such combinations were nonetheless
assayed, their effects are independently measured, and we evaluate the extrapolation against that
measurement in two ways: as the correlation between predicted and measured effect profiles, and as
the accuracy of the interaction classes assigned from the generated joint effect.

\subsubsection{Genetic interaction metrics}
\label{sec:gi-scoring}

We quantify the interaction type of each pair by Norman-style regression in decoded gene-expression
space. Writing the additive prediction as $\delta^{\mathrm{add}}_{ab} = \delta_a + \delta_b$, the
double effect is regressed on its two single effects,
\begin{align}
  \delta_{ab} \;\approx\; c_1\,\delta_a + c_2\,\delta_b,
  \label{eq:norman}
\end{align}
by Theil--Sen regression without intercept, giving coefficients $(c_1,c_2)$ and fitted double
$\hat{\delta}_{ab} = c_1\delta_a + c_2\delta_b$. Note that $(c_1,c_2)$ are post hoc regression
coefficients in gene space and are distinct from the model's parent coefficients $(c_a,c_b)$ of
\eqref{eq:parent-scale}. Table~\ref{tab:gi-metrics} lists the eleven scalar metrics derived from
$(\delta_a, \delta_b, \delta_{ab}, \hat{\delta}_{ab}, c_1, c_2)$. The same routine is applied to the
effect vectors of any method under comparison, so that only the effect vectors, and not the scoring
convention, differ between methods.

\begin{table}[!htbp]
\centering
\small
\begin{tabular}{lll}
\toprule
Metric & Symbol & Definition \\
\midrule
Fit magnitude            & $m$        & $\sqrt{c_1^2 + c_2^2}$ \\
Dominance                & $D$        & $\big|\log_{10}(|c_1|/|c_2|)\big|$ \\
Parent asymmetry         & $\Delta_{\!\rho}$ & $|\,\mathrm{dCor}(\delta_a,\delta_{ab}) - \mathrm{dCor}(\delta_b,\delta_{ab})\,|$ \\
Contribution balance     & $E$        & $\min / \max$ of the same two distance correlations \\
Single-effect similarity & $S$        & $\mathrm{dCor}(\delta_a, \delta_b)$ \\
Additive residual        & $R$        & $\|\delta_{ab} - \delta^{\mathrm{add}}_{ab}\| \,/\, \|\delta^{\mathrm{add}}_{ab}\|$ \\
Joint-to-additive ratio  & $\eta$     & $\|\delta_{ab}\| \,/\, \|\delta^{\mathrm{add}}_{ab}\|$ \\
Residual alignment       & $\alpha$   & $\cos\big(\delta_{ab} - \delta^{\mathrm{add}}_{ab},\; \delta_{ab}\big)$ \\
Fit quality              & $F$        & $\mathrm{dCor}(\delta_{ab}, \hat{\delta}_{ab})$ \\
Fitted-effect norm       & $\nu$      & $\|\hat{\delta}_{ab}\|$ \\
Fitted-to-joint ratio    & $\phi$     & $\|\hat{\delta}_{ab}\| \,/\, \|\delta_{ab}\|$ \\
\bottomrule
\end{tabular}
\caption{Genetic-interaction metrics. $\mathrm{dCor}$ denotes distance correlation.}
\label{tab:gi-metrics}
\end{table}

\subsubsection{Interaction classification}
\label{sec:gi-classification}

Each metric is $z$-scored across the pairs being classified. Every class is assigned a composite
score equal to the mean of its signed $z$-scored terms (Table~\ref{tab:gi-classes}), and each pair
receives the class with the highest score. Pairs with any non-finite metric are left unclassified.
The reference annotation distinguishes two synergy subtypes by whether the two single perturbations
induce similar or dissimilar phenotypes; we do not make this distinction and score synergy as one
class throughout, collapsing the reference subtypes accordingly.

\begin{table}[!htbp]
\centering
\small
\begin{tabular}{lp{7.2cm}l}
\toprule
Class & Biological interpretation & Composite score \\
\midrule
Approximately additive & Double effect explained by the sum of the singles & $-R - \eta$ \\[3pt]
Epistasis              & One parent dominates and masks the other & $+D + \Delta_{\!\rho} - E$ \\[3pt]
Redundant              & Parents act on a shared program; the double adds little beyond either alone & $+S - \eta$ \\[3pt]
Suppression            & Joint response damped relative to what both parents predict & $-\eta + \phi - m$ \\[3pt]
Neomorphic             & Double not explained by any weighting of the parents & $-F - \nu$ \\[3pt]
Potentiation           & Double exceeds the additive expectation; parents act on dissimilar programs & $+\eta + \alpha - S$ \\[3pt]
Synergy                & Double exceeds the additive expectation; parents act on similar programs & $+m + \eta + 2S$ \\
\bottomrule
\end{tabular}
\caption{Composite scores over $z$-scored metrics; each is averaged over its terms.}
\label{tab:gi-classes}
\end{table}
